\documentclass[10pt,a4paper]{article}
\usepackage[margin=0.58in]{geometry}
\usepackage[T1]{fontenc}
\usepackage{lmodern}
\usepackage{microtype}
\usepackage{enumitem}
\usepackage[hidelinks]{hyperref}
\usepackage{xcolor}
\usepackage{titlesec}
\usepackage{tabularx}
\usepackage{array}
\usepackage{parskip}
\definecolor{accent}{HTML}{1F1F1F}
\pagestyle{empty}
\setlength{\parindent}{0pt}
\setlength{\parskip}{0pt}
\setlist[itemize]{leftmargin=1.18em,itemsep=1.1pt,topsep=1.5pt,parsep=0pt,partopsep=0pt}
\titleformat{\section}{\large\bfseries\color{accent}}{}{0em}{}[\vspace{-0.40em}\rule{\textwidth}{0.32pt}]
\titlespacing*{\section}{0pt}{0.55em}{0.28em}
\newcommand{\entry}[4]{\begin{tabularx}{\textwidth}{@{}X>{\raggedleft\arraybackslash}p{0.28\textwidth}@{}}\textbf{#1} & #2 \\\textit{#3} & \textit{#4}\end{tabularx}}
\newcommand{\project}[4]{\begin{tabularx}{\textwidth}{@{}X>{\raggedleft\arraybackslash}p{0.23\textwidth}@{}}\textbf{#1} --- \textit{#2} & #3 \\\multicolumn{2}{@{}l@{}}{\footnotesize #4}\end{tabularx}}
\begin{document}
{\LARGE\bfseries Zhiheng Zhang}\\[-0.05em]
\href{mailto:zhiheng0mera@gmail.com}{zhiheng0mera@gmail.com}
\quad $\vert$ \quad
\href{mailto:zhihezhang@student.unimelb.edu.au}{zhihezhang@student.unimelb.edu.au}\\
\href{https://github.com/zhiheng-zhang-Mera}{github.com/zhiheng-zhang-Mera}
\hfill Melbourne, Australia
\vspace{0.30em}

\section{Research Profile}
Computer Science research engineer focused on mobile edge computing, distributed AI services and reliable multi-device execution.

\section{Education}
\entry{University of Melbourne}{Melbourne, Australia}{Master's degree in Computer Science}{Expected 2026}
\entry{University of British Columbia}{Canada}{Bachelor of Science in Computer Science}{Completed Summer 2024}

\section{Selected Research and Engineering Projects}
\project{Utopia / Digital City}{Multi-device personal-computing and research platform}{2026--present}{\href{https://github.com/zhiheng-zhang-Mera/Utopia}{github.com/zhiheng-zhang-Mera/Utopia}}
\begin{itemize}
  \item Built Android and Web control surfaces around a platform-neutral City protocol with authenticated device/control paths, explicit confirmation and truthful degraded-state handling.
  \item Current product evidence includes accepted V0.2/V0.3 boundaries and an accepted three-end mesh across two real Windows workers plus an Android control surface with strict target-device routing and cross-host review.
  \item The research-strengthening programme has accepted experiment registry, trace/provenance, repeated-scenario, fault-injection, and replay/ablation stages; artifact export remained under independent review at the 2026-10-06 snapshot.
  \item Next research direction: a \textbf{Personal Compute Fabric} for heterogeneous-device telemetry, placement/offloading, recovery and explainable cross-device execution; this is a planned extension, not claimed complete.
\end{itemize}

\project{DS-Hns}{Long-horizon autonomous software-work runtime}{2026--present}{\href{https://github.com/zhiheng-zhang-Mera/DS-Hns}{github.com/zhiheng-zhang-Mera/DS-Hns}}
\begin{itemize}
  \item Building long-horizon repository execution with durable task state, restart/recovery, hardware-aware concurrency and explicit terminal-state evidence.
  \item Uses failure isolation and recovery-oriented execution to study when a software agent has genuinely completed work rather than merely survived as a process.
  \item Supports research on coding-agent reliability, autonomous maintenance, AIOps and empirical software-engineering evaluation.
\end{itemize}

\project{Codex Boss}{Agent orchestration and evidence-based acceptance}{2026--present}{\href{https://github.com/zhiheng-zhang-Mera/Codex-Boss}{github.com/zhiheng-zhang-Mera/Codex-Boss}}
\begin{itemize}
  \item Developed a multi-agent workflow platform for planning, execution, critique and adjudication with durable state and evidence-aware acceptance.
  \item Treats unresolved failure cases and incomplete acceptance conditions as explicit evidence rather than reporting them as solved.
  \item Supports research on trustworthy agent orchestration, completion claims, auditability and human oversight.
\end{itemize}

\section{Research Interests}
Mobile edge computing; distributed systems; edge intelligence; resource-aware agent execution; digital-twin and multi-device systems.

\section{Technical and Research Strengths}
\begin{itemize}
  \item End-to-end research engineering: system design, implementation, testing, debugging, controlled experiments and evidence-driven iteration.
  \item Multi-device and long-running software systems with explicit lifecycle, recovery, routing, observability and acceptance boundaries.
  \item Git/GitHub-based reproducible development, regression testing, cross-host review and structured preservation of failure evidence.
\end{itemize}

\section{Current Research Direction}
I want to study how heterogeneous personal devices can form a measurable edge-computing substrate for agent workloads, with explicit placement, recovery and provenance.

\end{document}
